% The bus with the old device on it, and the three ways out.
\begin{tikzpicture}[font=\sffamily\scriptsize, x=1cm, y=1cm,
  tinynode/.style={busnode, text width=9mm, minimum height=5mm, inner sep=1.5pt,
               font=\sffamily\tiny}]

% ================================================================ the failing bus
\node[chlabel, anchor=west] at (-6.4,3.1) {the bus this chapter breaks on purpose};
\draw[busline] (-5.4,2.4) -- (5.4,2.4);
\node[busterm, anchor=east] at (-5.4,2.4) {};
\node[busterm, anchor=west] at (5.4,2.4) {};

\node[busnode, text width=24mm] (jn) at (-3.6,1.4) {the node};
\node[busmaster, text width=26mm] (ms) at (0,1.4) {the motion master};
\node[busoff, text width=26mm] (ol) at (3.8,1.4) {the classic hat, bought for this chapter};

\draw[busstub] (-3.6,2.4) -- (jn.north);
\draw[busstub] (0,2.4) -- (ms.north);
\draw[busstub] (3.8,2.4) -- (ol.north);

\node[buslabel, anchor=west, text=red!60!black] at (-5.3,0.5)
  {classic frames between all three: fine. Frames in the newer format: invalidated, every one of them};

% ================================================================ way out 1
\node[chlabel, anchor=north] at (-4.6,-0.2) {1. a purchasing rule};
\draw[busline] (-6.3,-1.3) -- (-2.9,-1.3);
\node[tinynode] (p1) at (-5.75,-2.0) {newer};
\node[tinynode] (p2) at (-4.6,-2.0) {tolerant};
\node[tinynode] (p3) at (-3.45,-2.0) {newer};
\draw[busstub] (-5.75,-1.3) -- (p1.north);
\draw[busstub] (-4.6,-1.3) -- (p2.north);
\draw[busstub] (-3.45,-1.3) -- (p3.north);

% ================================================================ way out 2
\node[chlabel, anchor=north] at (0,-0.2) {2. two buses};
\draw[busline] (-1.7,-1.1) -- (1.7,-1.1);
\node[tinynode] (q1) at (-0.85,-1.75) {newer};
\node[tinynode] (q2) at (0.85,-1.75) {newer};
\draw[busstub] (-0.85,-1.1) -- (q1.north);
\draw[busstub] (0.85,-1.1) -- (q2.north);
\draw[busline] (-1.7,-2.5) -- (1.7,-2.5);
\node[tinynode] (q3) at (-0.85,-3.15) {classic};
\node[tinynode] (q4) at (0.85,-3.15) {classic};
\draw[busstub] (-0.85,-2.5) -- (q3.north);
\draw[busstub] (0.85,-2.5) -- (q4.north);

% ================================================================ way out 3
\node[chlabel, anchor=north] at (4.6,-0.2) {3. a gateway};
\draw[busline] (2.9,-1.1) -- (6.3,-1.1);
\node[tinynode] (r1) at (5.7,-1.75) {newer};
\draw[busstub] (5.7,-1.1) -- (r1.north);
\node[busnode, text width=18mm, minimum height=8mm, font=\sffamily\tiny] (gw) at (3.9,-1.9)
  {this node, two controllers};
\draw[busstub] (3.9,-1.1) -- (gw.north);
\draw[busline] (2.9,-2.9) -- (6.3,-2.9);
\draw[busstub] (gw.south) -- (3.9,-2.9);
\node[tinynode] (r2) at (5.7,-3.55) {classic};
\draw[busstub] (5.7,-2.9) -- (r2.north);

% ================================================================ captions
\node[tlabel, anchor=north, text width=34mm, align=left] at (-4.6,-4.3)
  {Cheapest, and it constrains every device anyone adds to this machine from now on.};
\node[tlabel, anchor=north, text width=34mm, align=left] at (0,-4.3)
  {Complete, and it always works. It costs a second pair of wires, two more terminators and a second
   controller on somebody.};
\node[tlabel, anchor=north, text width=34mm, align=left] at (4.6,-4.3)
  {Firmware, latency, and it is not transparent: only payloads of eight bytes or fewer survive the
   downward journey.};

\node[umlnote, text width=112mm, anchor=north west] at (-6.4,-5.6)
  {All three are used in real machines and the first is the most common, because it is a decision
   made once at purchasing time rather than a thing to maintain. The third is the only one that lets
   an existing classic device stay exactly where it is, which is why it survives in retrofits.};
\end{tikzpicture}
